
\subsection{Free Congregation in a Free Church}

Source: “Folketbladet,” 1882, for August 10, 17, and 24, and September 7 and 14. Published the same year as a separate reprint under the title: “Free Congregation in a Free Church. Reply to the Thirty’s Declaration.” See the note above, p. 38. — Ed.

\begin{quote}
\emph{Note: This section is currently represented by an editorial summary. 
A full translation of the original text is planned and contributions are welcome.}
\end{quote}


This section (pp. 120–144) argues that the “Thirty’s Declaration” isn’t surprising because two distinct tendencies have long existed within the Conference, and the conflict is really about goals and methods rather than personal faults. The author refuses to trade personal accusations and instead aims to clarify what “the new direction” seeks for the Norwegian Free Church in America: Christ’s glory through the salvation and upbuilding of souls. Internal tensions are described as healthy and even necessary in a living fellowship—proof that the body is a community rather than a political party.

A central claim is that the Conference constitution is not the ultimate standard by which a movement should be judged. Constitutions can be revised and are meant to protect freedom, not to function as an infallible “touchstone.” The real measure must be Scripture. From Scripture and church history, the author defines the free church’s calling as building God’s kingdom through genuinely free congregations with real self-government—formed by Word and Sacraments, practicing mutual admonition and discipline, growing through love, and sending missionaries. Church history is read as a warning against two great corruptions: Roman Catholic clerical tyranny and state-church coercion. America’s freedom offers a chance to “begin again” on a right foundation—congregational freedom.

The essay then connects this calling to the “Open Declaration,” presenting it as a necessary blow against “church politics” and what the author calls a hollow Christianity—one satisfied with correct forms, external membership, and institutional power while the seriousness of faith and repentance fades. This hollow version of Christianity appears in both “papal” and “state” forms: one threatens punishment in the next world, the other in this one, while the free congregation points to Christ and warns that apart from him one withers. In the Norwegian-American context, the danger is named “Misconfinism,” described as teaching grace and absolution “without faith” and producing spiritual sleepiness. The “Open Declaration,” though sharp and not politically prudent, is defended as clearing rotten timber so that a true free-church home can be built.

A long portion elevates “Børnelærdommen” (the Catechism and Pontoppidan’s Explanation) as the church’s shared “mother’s milk” and the people’s spiritual inheritance. The author insists the church’s doctrinal unity and defense should rest on this common, congregationally-held foundation—not on elite clerical “theses” devised in pastors’ meetings and imposed as binding rules. The real problem is not honest theological discussion, but the use of theses and synodical decisions to bind consciences and silence congregations. Properly understood, the congregation has the right and duty to test teaching by God’s Word; pastors’ theories should be judged by the congregation’s confession, not the other way around.

On church union, the author supports “peace and unity on the ground of truth,” proposing three practical pillars: the congregation, the shared catechetical faith, and a common seminary. Schisms, the author argues, largely came from overvaluing a few pastors’ opinions and under-valuing congregations’ inherited Lutheran faith. True reconciliation requires restoring congregational authority, recognizing congregations as Lutheran when they stand on the confession, and using the shared catechism as the basis for cooperation.

Finally, the author defines “free congregation” spiritually and practically. Freedom is first Christ’s victory and the Spirit’s work, not a product of constitutions. Practically, freedom means the whole body participates in service—men, women, young, old—without reducing ministry to the pastor alone. Calling a pastor adds gifts; it should not transfer the congregation’s authority. “Free congregation in a free church” also means churchwide cooperation that genuinely represents congregations so decisions can be embraced with wholehearted “yes,” while still honoring conscience: those unable to join a work should be left in peace. The piece closes with a call for greater spiritual life, humility, and restraint in rhetoric, so the movement gains friends by stronger deeds and wiser speech.